La mAP50:95 disminuye de 0.4358 en cámaras conocidas a 0.2625 en cámaras nuevas.
Los intervalos del 95\,\% son $[0.4288,0.4785]$ y $[0.2538,0.3023]$,
respectivamente. La comparación con trece categorías comunes conserva una brecha
de 18.08 puntos porcentuales. Es consistente con dificultad de transferencia
entre cámaras, pero los intervalos no son una prueba causal del efecto del fondo
ni estiman variación entre semillas o nuevas regiones. El modelo fue seleccionado
en validación; los resultados de prueba se conservan como evaluación de ese ajuste.

Los objetos pequeños tienen mAP50:95 de 0.2792 en cis\_test y 0.1619 en trans\_test.
En cámaras nuevas, los medianos alcanzan 0.3220 y los grandes 0.2870; el rendimiento
no crece monótonamente con el área. La composición de categorías y escenas difiere
entre estratos, por lo que el tamaño no se aísla como una causa mediante estas tablas.
bird pasa de AP50:95 0.3334 a 0.0383 y rodent de 0.4269 a 0.0016, con 179 y 26
cajas de prueba trans, respectivamente. Estas caídas necesitan evaluación específica.
Fox obtiene AP cero con solo dos referencias cis y una trans, después de cinco cajas
de entrenamiento; badger tiene nueve cajas de entrenamiento. Esos soportes impiden
considerar la media global como evidencia de detección fiable de categorías raras.

Los doce casos Grad-CAM incluyen nueve mapas con variación espacial y tres planos.
En algunos aciertos se resaltan partes del animal; los falsos positivos también
resaltan vegetación, suelo o bordes de la entrada. Estas observaciones son exploratorias.
Los mapas planos de opossum no permiten localizar contribuciones positivas para
la propuesta y capa fijadas. No se reemplazaron los casos ni se modificó la técnica
congelada para obtener visualizaciones favorables. Los rótulos de horario no son
mediciones de iluminación, especialmente cuando una fotografía oscura se rotula day.
La galería no estima la frecuencia de estos patrones ni demuestra dependencia causal
del fondo. Para comprobarla harían falta intervenciones controladas sobre las imágenes.

La submuestra aceptada muestra menor AP en oclusión media y alta que en baja.
Sus intervalos del 95\,\% remuestrean las 7\,438 secuencias de trans\_test, con
200 réplicas por nivel, semilla 42 y modelo fijo. No capturan desacuerdo entre
revisores, notas faltantes, variación del entrenamiento ni cámaras futuras.
Los 18 casos excluidos son falta de respuesta humana: las proporciones de los
582 aceptados son descriptivas y no se reutilizan intervalos poblacionales ni
la garantía de planificación de los 600 originales. La composición de especies,
el soporte y el tratamiento crowd impiden atribuir toda diferencia a la oclusión.
La demostración de la cámara 7 muestra una agregación reproducible de fotografías;
sus 1\,420 fotos con detecciones y once categorías no son un censo de individuos.

La ampliación v2 incorpora solo negativos en ocho cámaras adicionales; no agrega animales de esas cámaras.
El fondo puede entonces ayudar a distinguir la procedencia de los negativos. La resolución original
de esa cohorte también difiere de las imágenes ECCV reducidas. El resultado se describe como
entrenamiento ampliado y no como reproducción estricta del train ECCV. Una comparación futura
necesitaría otra ejecución controlada, conservando los mismos tests y sin ajustar a sus resultados.
Un falso negativo de una especie amenazada puede invisibilizar su presencia en decisiones de manejo,
mientras que un falso positivo consume tiempo de revisión y puede crear registros de distribución incorrectos.
El umbral 0.25 de la demostración no se considera óptimo para todas las especies ni para una decisión de conservación.

Debe mantenerse revisión profesional de observaciones sensibles, documentar esfuerzo de muestreo
y validar en cámaras y hábitats del destino. Una tasa baja de detecciones también puede reflejar
fallos del sensor o cambio de iluminación. Los conteos de imágenes no estiman población ni individuos.
La separación entre cámaras ayuda a evaluar generalización, pero el entrenamiento en Estados Unidos
no justifica despliegue en Amazonía o Andes sin validación local.

Para registros sensibles se deben conservar fotografía, caja, confianza y decisión
del especialista. La etiqueta automática no constituye un registro confirmado.
La revisión debe considerar el costo de omisión, el de identificación incorrecta
y la capacidad del equipo de campo. Un umbral global no sustituye un protocolo
por especie ni el control del esfuerzo por cámara. Los soportes, exclusiones y
categorías sin referencia deben acompañar cada comparación: una AP indefinida no
demuestra ausencia biológica. La galería orienta nuevas muestras, pero no mide
frecuencia de errores. Las decisiones de manejo conservan responsabilidad humana
y deben admitir corrección cuando aparezcan evidencias adicionales.
